\documentclass[12pt]{article}
\usepackage[utf8]{inputenc}
\usepackage[T1]{fontenc}
\usepackage{hyperref}
\title{Sighu Language Resources for K-12 Curriculum Development}
\author{Compiled by Claude Sonnet 5.5 (claude-sonnet-5-5) for Andr\'as Kornai}
\date{}
\begin{document}
\maketitle

\section{Introduction}
Sighu (also Lesighu, Mississiou) is a Bantu language of Gabon, Guthrie code B.202. German Wikipedia places it between Koulamoutou and Lastoursville in Ogoou\'e-Lolo province. About 1,000 speakers are reported, based on a 1990 figure. English Wikipedia describes the language as undocumented. Glottolog classifies it in a different Bantu branch (Ndasaic/Samayic) from the Kele group that German Wikipedia names. Joshua Project notes that French is increasingly used in daily life. Search found no teaching materials of any kind.

\section{Primary Educational Resources}
No primers, textbooks, or school programmes were found.

\section{Linguistic and Reference Materials}
\begin{itemize}
\item English Wikipedia article on the Sighu language. \url{https://en.wikipedia.org/wiki/Sighu_language}
\item German Wikipedia article on Sighu. \url{https://de.wikipedia.org/wiki/Sighu}
\item ASJP (Automated Similarity Judgment Program) entry for Sighu. This is a word-list database, so it is a possible source of basic vocabulary. I did not inspect its contents. \url{https://asjp.clld.org/languages/SIGHU}
\item Glottolog languoid page for Sighu. \url{https://glottolog.org/resource/languoid/id/sigh1238.bigmap.html}
\end{itemize}

\section{Religious and Community Resources}
\begin{itemize}
\item Joshua Project profile of the Sighu people of Gabon. It is a mission-oriented demographic profile, not a language resource. \url{https://joshuaproject.net/people_groups/19288/GB}
\end{itemize}
No scripture translations or community organizations were found.

\section{Digital and Audio Resources}
No audio, video, corpus, or software resources were found.

\section{Curriculum Development Context}
No grammar, dictionary, or orthography was found in search results. Any Sighu-medium curriculum would first require fieldwork-based documentation, working with the speakers, who are mostly older adults. Until then, the language is best represented in school materials only through cultural and heritage content. Gabon's schooling is conducted in French. Local contacts and further literature (for example Gabonese linguistics at Omar Bongo University, or SIL Gabon) were not found in search and should be investigated separately.
\end{document}